\documentclass[notes,10pt, aspectratio=169]{beamer}

\usepackage{pgfpages}
% These slides also contain speaker notes. You can print just the slides,
% just the notes, or both, depending on the setting below. Comment out the want
% you want.
\setbeameroption{hide notes} % Only slide
%\setbeameroption{show only notes} % Only notes
%\setbeameroption{show notes on second screen=right} % Both

\usepackage{helvet}
\usepackage[default]{lato}
\usepackage{array}
\usepackage{eurosym}
\usepackage{setspace}
\usepackage{animate}
\usepackage{multicol}
\usepackage{breqn}
\usepackage{subcaption}
\usepackage{caption}
\captionsetup{font=scriptsize, justification=raggedright,singlelinecheck=false}
\usepackage[export]{adjustbox}
\usepackage[scale=2]{ccicons}
\usepackage{tikz}
\newenvironment{wideitemize}{\itemize\addtolength{\itemsep}{10pt}}{\enditemize}
\usepackage{natbib}
%\usepackage{enumitem}

\usepackage{tikz}
\usepackage{verbatim}
\setbeamertemplate{note page}{\pagecolor{yellow!5}\insertnote}
\usepackage{amsmath}
%\usepackage{mathpazo}
\usepackage{hyperref}
\usepackage{lipsum}
\usepackage{multimedia}
\usepackage{multirow}
\usepackage{graphicx}
\usepackage{dcolumn}
\usepackage{bbm}
\newcolumntype{d}[0]{D{.}{.}{5}}

\usepackage{changepage}
\usepackage{appendixnumberbeamer}
% \newcommand{\beginbackup}{
%   \newcounter{framenumbervorappendix}
%   \setcounter{framenumbervorappendix}{\value{framenumber}}
%   \setbeamertemplate{footline}
%   {
%      \leavevmode%
%      \hline
%      box{%
%       \begin{beamercolorbox}[wd=\paperwidth,ht=2.25ex,dp=1ex,right]{footlinecolor}%
% %         \insertframenumber  \hspace*{2ex} 
%       \end{beamercolorbox}}%
%      \vskip0pt%
%   }
%  }


\usepackage{minted}   % code prettyfier
\usepackage{xsavebox} % file-size-efficient saveboxes
\usepackage{animate}  % for animated scrolling
\usepackage{MnSymbol} % \triangle, \triangledown for scroll buttons
\usepackage{media9}   % buttons

%%%%%%%%%%%%%%%%%%%%%%%%%%%%%%%%%%%%%%%%%%%%%%%%%%%%%%%%%%%%%%%%%%%%%%%%%%%%%%
% \smoothScroll[<animate opts>]    % autoplay,loop,... (see: texdoc animate)
%              {<xsavebox id>}
%              {<viewport height>}
%              {<steps>}           % scrolling granularity
%              {<steps per sec>}   % while playing; >25 doesn't make sense
%%%%%%%%%%%%%%%%%%%%%%%%%%%%%%%%%%%%%%%%%%%%%%%%%%%%%%%%%%%%%%%%%%%%%%%%%%%%%%
\newcommand\smoothScroll[5][]{%
  \savebox\myMeasBox{\xusebox{#2}}%
  \edef\mywd{\the\wd\myMeasBox}%
  \edef\myht{\the\ht\myMeasBox}%
  \edef\mytht{\the\dimexpr\ht\myMeasBox+\dp\myMeasBox\relax}%
  \edef\portht{\the\dimexpr#3\relax}%
  \begin{animateinline}[#1,label={#2},width=\mywd,height=\portht]{#5}%
    \multiframe{#4}{%
      dRaiseLen=\the\dimexpr-\myht+\portht\relax+\the\dimexpr(\mytht-\portht)/%
                \numexpr#4-1\relax\relax
    }{%
      \begin{minipage}[b][\portht][b]{\mywd}%
        \raisebox{\dRaiseLen}[0pt][0pt]{\xusebox{#2}}%
      \end{minipage}%
    }%
  \end{animateinline}%
}
\newsavebox\myMeasBox % for measuring purposes
%%%%%%%%%%%%%%%%%%%%%%%%%%%%%%%%%%%%%%%%%%%%%%%%%%%%%%%%%%%%%%%%%%%%%%%%%%%%%%

%%%%%%%%%%%%%%%%%%%%%%%%%%%%%%%%%%%%%%%%%%%%%%%%%%%%%%%%%%%%%%%%%%%%%%%%%%%%%%
%% \topScrollButton{<xsavebox id>}
%%%%%%%%%%%%%%%%%%%%%%%%%%%%%%%%%%%%%%%%%%%%%%%%%%%%%%%%%%%%%%%%%%%%%%%%%%%%%%
\newcommand\topScrollButton[1]{%
  \mediabutton[
    jsaction={
      if(event.shift){anim['#1'].pause();anim['#1'].frameNum=0;}
      else try{anim['#1'].frameNum--}catch(e){}
    }
  ]{\fboxsep=0pt\framebox[\widthof{\xusebox{#1}}][c]{%
    \tiny\strut\raisebox{-0.2\height}{$\triangle\triangle\triangle$}}%
  }%
}  
%%%%%%%%%%%%%%%%%%%%%%%%%%%%%%%%%%%%%%%%%%%%%%%%%%%%%%%%%%%%%%%%%%%%%%%%%%%%%%

%%%%%%%%%%%%%%%%%%%%%%%%%%%%%%%%%%%%%%%%%%%%%%%%%%%%%%%%%%%%%%%%%%%%%%%%%%%%%%
%% \botScrollButton{<xsavebox id>}
%%%%%%%%%%%%%%%%%%%%%%%%%%%%%%%%%%%%%%%%%%%%%%%%%%%%%%%%%%%%%%%%%%%%%%%%%%%%%%
\newcommand\botScrollButton[1]{%
  \mediabutton[
    jsaction={
      if(event.shift){anim['#1'].pause();anim['#1'].frameNum=anim['#1'].numFrames-1;}
      else try{anim['#1'].frameNum++}catch(e){}
    }
  ]{\fboxsep=0pt\framebox[\widthof{\xusebox{#1}}][c]{%
    \tiny\strut\raisebox{0.1\height}{$\triangledown\triangledown\triangledown$}}%
  }%
}  
%%%%%%%%%%%%%%%%%%%%%%%%%%%%%%%%%%%%%%%%%%%%%%%%%%%%%%%%%%%%%%%%%%%%%%%%%%%%%%


 \setbeamertemplate{footline}{
    \hbox{%
    \begin{beamercolorbox}[wd=\paperwidth,ht=3ex,dp=1.5ex,leftskip=2ex,rightskip=2ex]{page footer}%
        %\usebeamerfont{title in head/foot}%
        %\insertshorttitle \hfill
        %\insertsection \hfill
        %\tableofcontents[currentsection] \hfill
        \insertsectionnavigationhorizontal{0.7\paperwidth}{}{} \hfill 
        \insertframenumber{} / \inserttotalframenumber
    \end{beamercolorbox}}%
}
\newcommand{\backupend}{
   \addtocounter{framenumbervorappendix}{-\value{framenumber}}
   \addtocounter{framenumber}{\value{framenumbervorappendix}} 
}


\usepackage{graphicx}
\usepackage[space]{grffile}
\usepackage{booktabs}

% These are my colors -- there are many like them, but these ones are mine.
\definecolor{blue}{RGB}{0,114,178}
\definecolor{red}{RGB}{213,94,0}
\definecolor{yellow}{RGB}{240,228,66}
\definecolor{green}{RGB}{0,158,115}

\setbeamercolor{section in head/foot}{fg=blue}

\useinnertheme[shadow]{rounded}
\setbeamercolor{block title example}{fg=black,bg=orange!50!white}
\setbeamercolor{block body example}{fg=black, bg=orange!30!white}
\setbeamercolor{block title alert}{fg=black,bg=green!50!white}
\setbeamercolor{block body alert}{fg=black, bg=green!30!white}
\setbeamercolor{block title}{use=structure,fg=structure.fg,bg=structure.fg!20!bg}
\setbeamercolor{block body}{parent=normal text,use=block title,bg=block title.bg!50!bg}

\hypersetup{
  colorlinks=false,
  linkbordercolor = {white},
  linkcolor = {blue}
}


%% I use a beige off white for my background
\definecolor{MyBackground}{RGB}{255,253,218}

%% Uncomment this if you want to change the background color to something else
%\setbeamercolor{background canvas}{bg=MyBackground}

%% Change the bg color to adjust your transition slide background color!
\newenvironment{transitionframe}{
  \setbeamercolor{background canvas}{bg=yellow}
  \begin{frame}}{
    \end{frame}
}

\setbeamercolor{frametitle}{fg=blue}
\setbeamercolor{title}{fg=black}
%\setbeamertemplate{footline}[frame number]
\setbeamertemplate{navigation symbols}{} 
\setbeamertemplate{itemize items}{-}
\setbeamercolor{itemize item}{fg=blue}
\setbeamercolor{itemize subitem}{fg=blue}
\setbeamercolor{enumerate item}{fg=blue}
\setbeamercolor{enumerate subitem}{fg=blue}
\setbeamercolor{button}{bg=white,fg=blue,}



% If you like road maps, rather than having clutter at the top, have a roadmap show up at the end of each section 
% (and after your introduction)
% Uncomment this is if you want the roadmap!
% \AtBeginSection[]
% {
%    \begin{frame}
%        \frametitle{Roadmap of Talk}
%        \tableofcontents[currentsection]
%    \end{frame}
% }
\setbeamercolor{section in toc}{fg=blue}
\setbeamercolor{subsection in toc}{fg=red}
%\setbeamersize{text margin left=1em,text margin right=1em} 
\setbeamersize{text margin left=0.7cm,text margin right=0.7cm} 


\title[]{\textcolor{blue}{Practice 5: OVB's Hide and Seek}}
\author{25117 - Econometrics}
\institute[shortinst]{Universitat Pompeu Fabra}

\date{November 28th, 2023}

\begin{document}


% Title Slide

\begin{frame}[plain]
\maketitle
\end{frame}

\begin{frame}{What we learned in the last lesson}
\begin{itemize}
\item When $Y$ is a binary variable, the population regression function shows the probability that $Y = 1$ given the value of the regressors, $X_1, X_2,\dots, X_k$.
\item The linear multiple regression model is called the linear probability model when $Y$ is a binary variable because the probability that $Y = 1$ is a linear function of the regressors.
\item Probit and logit regression models are nonlinear regression models used when Y is a binary variable. Unlike the linear probability model, probit and logit regressions ensure that the predicted probability that $Y = 1$ is between 0 and 1 for all values of $X$.
\item Probit regression uses the standard normal cumulative distribution function. Logit regression uses the logistic cumulative distribution function. Logit and probit coefficients are estimated by maximum likelihood.
\item The values of coefficients in probit and logit regressions are not easy to interpret. Changes in the probability that $Y = 1$ associated with changes in one or more of the $X$’s depends on the initial value of $X$.
\item Hypothesis tests on coefficients in the linear probability, logit, and probit models are performed using the usual t- and F-statistics.
\end{itemize}
\end{frame}


\begin{frame}{Practice Setting}
\begin{block}{Your Research Question}
    What is the elasticity-price of cigarettes?
\end{block}\vspace{.5cm}
\textit{Why are we interested in knowing the elasticity of demand for cigarettes?}\\[1em]

\begin{itemize}
    %\item \textbf{Theory of optimal taxation --- } The optimal tax rate is inversely related to the price elasticity: the greater the elasticity, the less quantity is affected by a given percentage tax, so the smaller is the change in consumption and deadweight loss.\\[1em]
    \item \textbf{Non-monetary externalities ---} role for government intervention to discourage smoking health effects of second-hand smoke? \\[1em]
    \item \textbf{Monetary externalities ---} Pressure on public health system
\end{itemize}
\vfill
\textit{By the way... what is an externality?}

\end{frame}


\begin{frame}{Threats to Internal Validity}

\begin{exampleblock}{Question}
    What is internal validity? How does it differ from external validity?
\end{exampleblock}

Remember... There are 5 threats to the internal validity of regression studies:
\begin{itemize}
    \item Omitted variable bias
    \item Wrong functional form
    \item Errors-in-variables bias
    \item Sample selection bias
    \item Simultaneous causality bias
\end{itemize}\\[1em]
All of these imply that $\mathbb{E}(u_i\mid X_1, X_2, \dots, X_k)\neq 0$ (or that the conditional mean independence fails) -- in which case OLS is biased and inconsistent!


\begin{exampleblock}{Question}
    What can we do to tackle each one of them?
\end{exampleblock}

\end{frame}



\begin{frame}{Panel Dataset}
\begin{exampleblock}{Question}
    What is a panel dataset? Why is it useful? A cross-sectional dataset?
\end{exampleblock}
\vspace{.5cm}
\textbf{Your panel:}
\begin{itemize}
    \item 48 continental U S states, 1985–1995
    \item Annual cigarette consumption
    \item Average prices paid by end consumer (including tax)
    \item Personal income
    \item Tax rates (cigarette-specific and general statewide sales tax rates)
\end{itemize}
\begin{exampleblock}{Question}
    What would be your main empirical strategy to answer the research question at hand?
\end{exampleblock}

\end{frame}


\begin{frame}{Estimation Strategy}
We need to combine \underline{both} IV and panel data methods! \\[.25em]
\begin{itemize}
    \item[$\longrightarrow$] We need unit fixed effects which control for unobserved state-level characteristics that affect the demand for cigarettes and the tax rate, as long as those characteristics don’t vary over time.
\end{itemize}
\begin{equation*}
    log(Q^{cigarettes}_{it}) = \alpha_i + \beta_1 log(P^{cigarettes}_{it}) + \beta_2 log(Income_{it}) + u_{it}
\end{equation*}
\begin{exampleblock}{Question}
    How should you interpret $\beta_1$? Why controlling for $Income$ ? 
\end{exampleblock}

\begin{itemize}
   \item[$\longrightarrow$] We need IV methods to handle the simultaneous causality bias that arises from the interaction of supply and demand.\\[1em]
\end{itemize}


\begin{exampleblock}{Question}
    What is simultaneous causality bias and how is it related to the problem at hand?
\end{exampleblock}

\end{frame}


\begin{frame}{Long Term Effects}
One way to account for state fixed effects is to model long-term effects (\textit{why?}), e.g., consider 10-year changes, between 1985 and 1995\\[.25em]

Rewrite the regression in “changes” form:
\begin{eqnarray*}
    &log(Q^{cigarettes}_{i,1995})\text{--}log(Q^{cigarettes}_{i,1985}) &= \beta_1 [log(P^{cigarettes}_{i,  1995}\text{--}log(P^{cigarettes}_{i,  1985})]\\
    & {} &+ \beta_2 [log(Income_{i,  1995})\text{--}log(Income_{i,  1985})]\\
    & {} &+ [u_{i,  1995}-u_{i,  1985}]
\end{eqnarray*}
Then estimate the demand elasticity by 2SLS using 10-year changes in the instrumental variables

\begin{exampleblock}{Question}
    How would you write the corresponding short-run specification, while accounting for state fixed effects, $\alpha_i$ ? 
\end{exampleblock}


\end{frame}


\begin{frame}{Long Term Effects}
\begin{table}[]
\begin{tabular}{l|ccc|}
\cline{2-4}
                                                                                                                               & \multicolumn{3}{c|}{ {$log(Q^{cigarettes}_{i,1995})$--$log(Q^{cigarettes}_{i,1985})$}} \\
                                                                                                                               & \multicolumn{3}{c|}{ } \\
                                                                                                                               & (1)                                        & (2)                                                     & (3)                                     \\ \hline
\multicolumn{1}{|l|}{$log(P^{cigarettes}_{i,  1995}$ -- $log(P^{cigarettes}_{i,  1985})$} & -0.094                                     & -1.34                                                   & -1.20                                   \\
\multicolumn{1}{|l|}{}                                                                                                         & (0.21)                                     & (0.23)                                                  & (0.20)                                  \\
\multicolumn{1}{|l|}{$log(Income_{i,  1995})$ -- $log(Income_{i,  1985})$}                                                      & 0.53                                       & 0.43                                                    & 0.46                                    \\
\multicolumn{1}{|l|}{}                                                                                                         & (0.34)                                     & (0.30)                                                  & (0.31)                                  \\
\multicolumn{1}{|l|}{Constant}                                                                                                 & -0.12                                      & -0.02                                                   & -0.05                                   \\
\multicolumn{1}{|l|}{}                                                                                                         & (0.07)                                     & (0.07)                                                  & (0.06)                                  \\ \hline
\multicolumn{1}{|l|}{Instrumental Variable(s)}                                                                                 & Sales tax                                  & Cigarette tax                                  & Both                                    \\
\multicolumn{1}{|l|}{First-stage-F-statistic}                                                                                  & 33.7                                       & 107.2                                                   & 88.6                                    \\
\multicolumn{1}{|l|}{J-Statistic}                                                                                              & --                                         & --                                                      & 4.93                                    \\
\multicolumn{1}{|l|}{}                                                                                      & --                                         & --                                                    & [0.026]                                 \\ \hline
\end{tabular}
\end{table}
{\footnotesize Robust standard errors are reported in parenthesis. p-values are reported in brackets.}
\end{frame}

\begin{frame}{Pop-up Quiz}
\begin{exampleblock}{Question(s)}
    \begin{itemize}
        \item[1.] Interpret the coefficients of interest
        \item[2.] Do they seem large? Small?
        \item[3.] How do we interpret the coefficients of the LR difference in income? What does it suggest?
        \item[4.] What are the two conditions required for a valid instrument?
        \item[5.] How to check for each condition?
        \item[6.] What is a J-test? 
        \item[7.] How do we interpret the J-statistic?
        \item[8.] Here, we are looking at long-run (ten-year change) elasticities. What would you expect a short-run (one-year change) elasticity to be --- more or less elastic?
        \item[9.] What are the remaining threats to internal validity?
        \item[10.] Discuss the external validity of the results
    \end{itemize}
\end{exampleblock}

\end{frame}


\begin{frame}{}
\begin{center}
    {\Large Good Luck for the Exam!}
\end{center}
\end{frame}


\end{document}